\section{The Synthesis of All Traditions}

\begin{quotex}
It is true that many vestiges of a forgotten past are coming out of the earth in our age, and perhaps not without reason.
\flright{\textsc{Rene Guenon}, \emph{Traditional Forms and Cosmic Cycles}}
\end{quotex}


Guenon wrote that well before the discoveries of the Dead Sea Scrolls and the Nag Hammadi Library. The Tibetans have a similar view and believe that sacred texts may be hidden for hundreds of years and then found at the opportune time. We here have noted the publication of texts that point to the likely existence of a continuous esoteric tradition in the West. Even \textbf{Dante}'s Divine Comedy can be understood in that vein. As Guenon points out, the inner meaning of that poem was hidden for six hundred years until some poets in the 19\textsuperscript{th} century began to suspect there was a deep metaphysical teaching in the Comedy. It seems, however, that little has been done in the meantime to build on that insight, at least until we encounter the unique meditations of Guido De Giorgio.

\textbf{Guido De Giorgio} was born on 3 October 1890 at San Lupo and died a natural death on 27 December 1957. After graduating with a degree in philosophy from the University of Naples, he emigrated to Tunisia to teach. While there, he was initiated into the Sufi order led by
\textbf{Sheik Mohammad Kheireddine}. He returned to Italy at the outbreak of the WW I. After that, he spent quite some time in Paris where he developed a close relationship with Rene Guenon. He returned to the Italian Alps and married his second wife who had done her university thesis on the Vedanta. One of his sons died a hero's death in Ethiopia; it was a loss, since he left behind a journal that displayed his father's influence. His daughter became a nun and his other son was named ``Renato'' (Italian for Rene) in honor of M. Guenon.

In Italy, De Giorgio collaborated with \textbf{Julius Evola} in the Ur group and then in Evola's journal \emph{La Torre}; he also contributed some articles to \emph{Diorama filosofico}, a philosophical journal also directed by Evola. Evola described him in these terms:

\begin{quotex}
He was a type of initiate in a primitive and chaotic state, he had both lived among the Arabs and known Guenon, who held De Giorgio in high esteem. He possessed an exceptional culture, knew several languages, but his temperament was unstable. His refusal of the modern world was such that he lived secluded in the mountains that he felt as his natural milieu and, finally, in an abandoned rectory, living practically on nothing, from some tutoring that he did, and suffering physically every time he was obliged to make contact with civilized and town life.

His influence on me, which owes nothing to his books, since none of them had been published yet, but which was exercised by his upsetting and aggressive letters, sprinkled with illuminations, and confusions, was tied to his way of dramatizing and stimulating his conception of Tradition. Guenon's presentation was, because of his personal equation, very formal and intellectual traits. But De Giorgio united to that a tendency of absolutism which, naturally, found in me congenital territory. The rare pieces of his that I published, sometimes extracts of his letters against his will, are perhaps the only ones that remain, unfortunately. I was in contact with De Giorgio (whom I met twice in the Alps) during the brief period of my review La Torre. Otherwise, in more recent times, we have become distant from each other, because of his drift into a type of ``Vedantizing'' Christianity.
\end{quotex}


Nevertheless, that is still too restrictive, because he could equally be accused of ``paganizing'' Christianity, or more especially, ``Sufi-izing'' Christianity. De Giorgio, along with \textbf{Arturo Reghini} and Julius Evola, was concerned with the project of the restoration of the Roman Tradition. \textbf{Piero di Vona} describes this:

\begin{quotationx}
Regarding the project of the restoration of the Roman tradition and with it, Traditional Europe, which is owed to De Giorgio, the best Italian disciple of Guenon, we can only mention a few things because of the great complexity and intrinsic difficulty of his work which would require a deep and independent study. Whoever wants to understand De Giorgio must have a deep knowledge of Guenon. Moreover, we have to warn any possible reader of De Giorgio's style of thought and way of writing. The inspired, wild, and almost raving tone, in some of his pages, masks a very deep understanding of the ideas and symbols of contemporary, metaphysical Traditionalism.

De Giorgio's vision of the transcendent truths of the traditional order is very complex. In it, Neoplatonic, Christian, Hindu, and Islamic themes join and intersect each other. His conception of the supreme principle as silence, and the three degrees of the ineffable, should be studied in relationship to Neoplatonism. Some profound Christian motifs are present.

The non-dual understanding of the supreme is the aspect of De Giorgio's conception that is tied to the well-known Hindu idea, which he shares in common with Guenon. But the Islamic perspective remains preeminent and dominant over all the indicated themes. But it is certainly important to note that, through Guenon, the first appearance in Italy of the Islamic vision of the absolute was made with De Giorgio, from which he began to build his Roman traditionalism. That is so much more interesting given that De Giorgio had never converted to Islam.
\end{quotationx}


There is no point to continue with this at this time, as it will become clearer as more of his essays appear in Gornahoor. Yet, things cannot be so simple with De Giorgio, since he also ties Tradition to the two Romes: the ancient pagan, and the medieval Catholic. To return to the beginning, ultimately it is Dante who brings it all together. On the one hand, Dante was influenced by the writings of the Sufi Ibn Al-Arabi. On the other, through the pagan \textbf{Virgil}, and ultimately \textbf{Aeneas}, the Indo-European pagan tradition enters into the Tradition of the second Rome. De Giorgio goes even further: Dante, for him, brings together the two dominant spiritual strains in the Middle Ages, those of the Dominicans and the Franciscans, or, to say it again, the Thomist and the Augustinian paths. As De Giorgio writes:

\begin{quotex}
The Dominicans and the Franciscans have marked out the two greatest ways of the realization of the divine.
\end{quotex}


Regarding his views on Dante, Vona summarizes De Giorgio:

\begin{quotex}
There is in Rome the living synthesis of the ancient pagan tradition and of the new and Christian one. Dante is the one who not only manifested it, but directly transmitted in the sacred poem and its realizing and initiatory values and meanings. One can reach the most profound parts of Catholicism from Dante, which are related to metaphysical realization by means of contemplation.
\end{quotex}

References:

\textsc{Baillet, Phillipe}, \textit{Guido De Giorgio, le voyant solitaire} from \textit{L'Instant et l'Eternité}

\textsc{Di Vona, Piero}, \textit{Evola Guenon De Giorgio}

\textsc{Pintore, Franco}, \textit{Preface to La Tradizione Romana}


\flrightit{Posted on 2013-02-15 by Cologero}

\begin{center}* * *\end{center}

\begin{footnotesize}
\begin{sffamily}
\texttt{Janus on 2013-02-15 at 01:24 said:}

Fascinating. It begs the question of whether a restored Tradition in the West will be Christian in nature, or whether it will take a form which draws from many other manifestations of Tradition around the world. Perhaps those who will become the guides must seek out strange lands and traditions to find what is so familiar in their own they become blind to it, before they can ``know the place for the first time.'' (Eliot).

\hfill

\texttt{Mercurius on 2013-02-15 at 15:07 said:}

Quite interesting and similarly, Theophilus Schweighardt begins his Rosy Cross treatise ``Mirror of Wisdom'' explaining that the whole teaching is being presented ``clearly'' and ``before your eyes''. Then, after pointing out that the ``Holy Divine Scripture'' comprises the bedrock of the doctrine, he refers readers to the Augustinian Thomas Kempis for practical instruction:

``More I cannot say unto thee of this, but if thou desirest more information concerning this fundament and preparatory work, thou shalt find more thereon in the aforementioned little books of Thomas a Kempis, for the author in the same book does nothing else but teach thee to practise this work rightly and well, and so it may be called his golden writing, well and truly a fount and origin of the Rhodo-staurotic teaching. Hoc de priori''.

\hfill

\texttt{Cologero on 2013-02-19 at 07:37 said:}

Perhaps we could distinguish the main points of the two schools in this manner. For the Thomist, all knowledge begins in sense perception. For the Augustinian, it begins in the interiority of man.


\hfill

\end{sffamily}
\end{footnotesize}

